\subsection{Proof by Mathematical Induction}
\label{subsections:techniques:universal:induction}



\paragraph{Description.}
A proof by induction proves a mathematical statement by showing that it %
holds for a first case and that each case implies the next. %
Precisely, it proves a universal statement, %
$\forall n \in S$ with $n \ge n_0$, $P(n)$, %
by showing that $P(n_0)$ is true and that $\forall k \in S$ with %
$k \ge n_0$, if $P(k)$, then $P(k + 1)$. %
Generally, a proof by induction mirrors the following format:
\begin{enumerate}{
    \item \textbf{Setup.} %
        Define the statement $P(n)$ and identify the smallest case $n_0$.
    \item \textbf{Base Case.} %
        Let $n \in S$ such that $n = n_0$ and show that $P(n_0)$ is true.
    \item \textbf{Inductive Hypothesis.} %
        Let $k \in S$ such that $k \ge n_0$ and $P(k)$ is true.
    \item \textbf{Inductive Step.} %
        Show that $P(k + 1)$ is true by using the inductive hypothesis.
    \item \textbf{Conclusion.} %
        Conclude the proof by acknowledging the principle of mathematical %
        induction, optionally by summarizing what has been found.
}
\end{enumerate}



\paragraph{Advice.}
\begin{itemize}{
    \item Use this proof when $P(k + 1)$ can be related to $P(k)$. %
        A summation sign is a common indication of induction.
    \item The inductive hypothesis always follows the same format, %
        copying $P(n)$ with $n$ replaced by $k$. %
        This is a common source of confusion, as $P(k)$ is assumed %
        rather than proven.
    \item In the inductive step, relate $P(k + 1)$ back to $P(k)$ to use %
        the inductive hypothesis, without assuming $P(k + 1)$. %
        Optionally, list what is to be proven before doing so.
}
\end{itemize}



\paragraph{Example.}
\begin{exercise}{
    Prove that $\forall n \in \Z$ with $n \ge 1$, %
    $\displaystyle\sum_{i=1}^{n} i = \frac{n(n + 1)}{2}$.
}
\end{exercise}
\begin{proof}{
    This statement can be proven by induction.
    \begin{annotatedpanel}[Setup.]{
        Let $P(n)$ be the statement that
        \[
            \sum_{i=1}^{n} i = \frac{n(n + 1)}{2}.
        \]
        Since $n \ge 1$, the smallest case is $n_0 = 1$.
    }{
        \textbf{Step 1.} Define $P(n)$ and identify the smallest case %
        $n_0$. %
        Here, $S = \Z$ and $n_0 = 1$, as the exercise requires $n \ge 1$.
    }
    \end{annotatedpanel}
    \begin{annotatedpanel}[Base Case.]{
        Let $n \in \Z$ such that $n = 1$. %
        Then
        \begin{align*}{
            \sum_{i=1}^{n} i &= \sum_{i=1}^{1} i = 1,\\
            \frac{n(n + 1)}{2} &= \frac{1(1 + 1)}{2}
                = \frac{1(2)}{2} = \frac{2}{2} = 1.
        }
        \end{align*}
        Thus, $1 = 1$, so $P(1)$ is true.
    }{
        \textbf{Step 2.} Let $n \in S$ such that $n = n_0$ and show that %
        $P(n_0)$ is true. %
        Here, $n_0 = 1$ is the smallest possible case based on $n \ge 1$.
    }
    \end{annotatedpanel}
    \begin{annotatedpanel}[Inductive Hypothesis.]{
        Let $k \in \Z$ such that $k \ge 1$ and $P(k)$ is true.
    }{
        \textbf{Step 3.} Let $k \in S$ such that $k \ge n_0$ and $P(k)$ %
        is true. %
        Here, $S = \Z$, $k$ is at least the largest base case, and $P(k)$ is %
        $P(n)$ copied with $n$ replaced by $k$.
    }
    \end{annotatedpanel}
    \begin{annotatedpanel}[Inductive Step.]{
        Consider $P(k + 1)$. %
        We want to prove that
        \[
            \sum_{i=1}^{k+1} i = \frac{(k + 1)((k + 1) + 1)}{2}.
        \]
        Note that
        \[
            \sum_{i=1}^{k+1} i = (k + 1) + \sum_{i=1}^{k} i.
        \]
        By the inductive hypothesis,
        \begin{align*}{
            \sum_{i=1}^{k+1} i &= (k + 1) + \frac{k(k + 1)}{2}\\
                &= (k + 1) + \frac{k^2 + k}{2}\\
                &= \frac{2k + 2}{2} + \frac{k^2 + k}{2}\\
                &= \frac{k^2 + 3k + 2}{2}\\
                &= \frac{(k + 1)(k + 2)}{2}.
        }
        \end{align*}
        Thus, $P(k + 1)$ is true.
    }{
        \textbf{Step 4.} Show that $P(k + 1)$ is true by using the %
        inductive hypothesis. %
        Here, the sum splits into $k + 1$ and the sum up to $k$, %
        which the inductive hypothesis replaces. %
        The line beginning ``We want to prove'' is optional.
    }
    \end{annotatedpanel}
    \begin{annotatedpanel}[Conclusion.]{
        Thus, $\forall n \in \Z$ with $n \ge 1$, %
        $\displaystyle\sum_{i=1}^{n} i = \frac{n(n + 1)}{2}$. %
        Therefore, the statement holds by the principle of mathematical %
        induction. \qedhere
    }{
        \textbf{Step 5.} Conclude the proof by acknowledging the principle %
        of mathematical induction. %
        Here, the summary states what has been found: %
        $P(1)$ is true, and $P(k + 1)$ is true whenever $P(k)$ is true.
    }
    \end{annotatedpanel}
}
\end{proof}